\subsection{Prospective Validation of the First-Generation Reference Pathway}

The immediate research priority is not expansion of the proposed system, but
prospective validation of the first-generation reference pathway identified in
Section~6. The central question is whether an EEG-guided, stage-aware,
closed-loop TMR system can reproduce the physiological control and measurable
behavioral effect required for the translational argument to remain valid.
Validation should therefore address two related but analytically distinct
problems: whether the intervention is executed as intended and whether correctly
executed intervention produces the expected memory effect.

\noindent\textbf{Intervention fidelity.}

The first requirement is to establish that closed-loop operation can be
reconstructed objectively from synchronized physiological and intervention
records. Validation should determine whether usable EEG can be maintained
during overnight recording, whether state estimates are generated causally,
whether low-quality or ambiguous intervals lead to abstention, and whether
actual cue onset occurs while the physiological evidence supporting the
decision remains valid. Post-cue monitoring should determine whether
stimulation is followed by arousal, stage transition, signal degradation, or
other changes relevant to subsequent cueing.

These outcomes should be evaluated at the intervention level rather than only
through whole-night summary statistics. For each potential cueing opportunity,
the system should permit reconstruction of the physiological evidence available
at the time, the resulting control decision, whether a cue was actually
delivered, and the physiological state surrounding cue onset. Such
intervention-level traceability would distinguish failure of sensing or control
from absence of a biological TMR effect.

No universal thresholds for state confidence, signal quality, latency, cue
spacing, or arousal detection are established by the present synthesis.
Initial studies should therefore estimate these operating ranges empirically
rather than embedding literature-derived values as fixed specifications. The
objective is not to maximize stimulation frequency, but to identify operating
conditions under which intervention decisions remain physiologically
interpretable.

\noindent\textbf{Behavioral efficacy under verified closed-loop operation.}

Behavioral validation should use controlled paradigms in which pre-sleep
learning, cue--memory associations, cue exposure, physiological state at
stimulation, and post-sleep retention can be linked at the item or trial level.
Cued and uncued material should be appropriately matched so that any observed
difference can be distinguished from ordinary sleep-dependent consolidation or
general practice effects.

The principal question is not merely whether mean memory performance improves,
but whether a reproducible TMR effect is observed when the system demonstrably
executes the intended protocol. Given the modest and heterogeneous effect
reported in the broader literature \citep{hu2020promoting}, a single positive
experiment would provide insufficient evidence for reference-system validation.
Replication across participants, recording nights, and preferably independent
experimental samples should establish whether the behavioral effect persists
when intervention fidelity is controlled.

Technical and behavioral outcomes need not be evaluated in completely separate
studies, but behavioral efficacy should not be interpreted independently of
technical fidelity. A null behavioral result obtained during unreliable state
estimation or mistimed cue delivery would not constitute a clean test of TMR.
Conversely, repeated null effects obtained despite verified physiological
operation would challenge the assumption that the selected reference pathway
preserves the intervention conditions necessary for measurable benefit.

This distinction creates an explicit falsification criterion for the
first-generation hypothesis. If reliable EEG acquisition, causal
state-dependent control, and disturbance-limited cue delivery can be
demonstrated but replicated behavioral benefit cannot, development should not
progress automatically toward greater algorithmic complexity. The scientific
premise, task selection, cueing protocol, or practical value of the pathway
would instead require re-evaluation.


\subsection{Identifying the Physiological Variables That Improve Control}

Once the reference pathway can be evaluated reproducibly, the next question is
which additional physiological information actually improves intervention
outcomes. Future development should distinguish variables that are merely
measurable from variables that provide causal or predictive value for deciding
when and how to stimulate.

\noindent\textbf{Adaptive control as a testable hypothesis.}

Individual and night-to-night differences in sleep architecture, sensory
threshold, pre-sleep memory strength, cue history, and response to previous
stimulation provide plausible inputs for adaptive control. However, adding
adaptive variables increases model flexibility and creates opportunities for
overfitting, unstable decision policies, and post-hoc explanations of
successful or unsuccessful trials.

Each adaptive feature should therefore be introduced through prospective
comparison with a simpler reference policy. A variable should influence cue
eligibility, intensity, spacing, repetition, or temporary suspension only when
its inclusion improves at least one predefined outcome: behavioral efficacy,
preservation of sleep, intervention precision, or reliability of abstention.
Prediction of outcome within the same dataset is not sufficient evidence that a
variable should control future stimulation.

Where personalization is investigated, validation should distinguish
population-level predictive effects from individual-specific adaptation.
Apparent subject-specific patterns may reflect noise when only a small number of
nights or cueing events are available. Adaptive policies should therefore be
evaluated prospectively and, where feasible, against fixed-policy or
cross-validated baselines before being interpreted as genuine personalization.

\noindent\textbf{Stage-aware versus event- and phase-aware targeting.}

A separate question concerns the temporal resolution at which TMR should
operate. The rationale for finer targeting is supported not only by the
association of memory processing with slow oscillations and sleep spindles, but
also by closed-loop experiments demonstrating phase-dependent behavioral and
electrophysiological effects
\citep{staresina2015nesting,cairney2018spindle,antony2018spindle,
shimizu2018closedloop,ngo2022shaping,nicolas2025phase}.
These findings establish phase-aware TMR as an experimentally supported
alternative rather than a purely theoretical extension. They do not yet
establish that phase-aware stimulation provides a robust and generalizable
incremental advantage over well-controlled stage-aware cueing across memory
domains and implementation conditions.

This question should therefore be tested directly rather than resolved through
mechanistic plausibility. Stage-aware and temporally refined policies should be
compared under matched learning tasks, sensing conditions, cue exposure, and
outcome measures. Evaluation should include behavioral effect size, trigger
precision, missed opportunities, end-to-end timing error, artifact
susceptibility, sleep disruption, and reproducibility across participants and
nights.

Fine temporal control should be adopted only if its additional physiological
specificity translates into a reproducible advantage that exceeds the cost of
greater detection and synchronization complexity. A failure to demonstrate such
an advantage would itself be informative: it would support retaining
stage-aware control as the more parsimonious intervention strategy rather than
treating increased temporal precision as inherently superior.


\subsection{Comparative Sensing and Sensor-Reduction Experiments}

The EEG-first decision made in Section~6 should ultimately be tested rather than
preserved by assumption. Once a sufficiently characterized EEG-guided reference
configuration exists, it can provide the experimental basis for determining how
much physiological sensing is actually necessary for effective TMR.

\noindent\textbf{Hybrid sensing and complementary physiological information.}

Initial comparative studies should acquire EEG and peripheral signals
synchronously rather than replacing EEG immediately. Candidate peripheral
signals include PPG, accelerometry, respiratory information, and other measures
that may contribute information about movement, arousal, sleep stability, or
signal reliability.

The relevant question is not whether peripheral sensors improve agreement with
conventional sleep-stage labels. Their value should be measured by whether they
improve intervention decisions. Examples include identifying intervals in which
EEG is unreliable, detecting physiological disturbance following stimulation,
supporting abstention decisions, or reducing uncertainty about state
transitions.

A hybrid configuration would therefore be justified when the additional signal
provides measurable information beyond the EEG reference that improves
intervention reliability, behavioral outcome, or preservation of sleep.
Conversely, a sensor should not be retained merely because its data can be
acquired.

\noindent\textbf{Reduced EEG and peripheral-dominant control.}

Sensor reduction should be evaluated prospectively as an intervention problem
rather than solely as a sleep-staging problem. Reduced-channel EEG,
EEG--peripheral hybrid systems, and eventually PPG-dominant controllers should
be compared with the reference configuration using synchronized recordings and,
where possible, identical cueing opportunities.

Relevant endpoints include agreement in cue-eligibility decisions, false
authorization of stimulation, abstention during uncertain intervals,
sensing-to-cue timing, detection of post-cue disturbance, preservation of sleep
continuity, and behavioral TMR outcome. A reduced-sensor system that reproduces
retrospective stage labels but systematically changes which cues are delivered
has not demonstrated equivalence at the intervention level.

Peripheral control should therefore not be required to reproduce every EEG
feature. The more meaningful question is whether the reduction in neural
observability produces an unacceptable loss of intervention validity.
A peripheral or hybrid pathway could justify progression if it preserves
behavioral efficacy and physiological safety within prospectively defined
bounds while materially reducing measurement burden.

This creates a testable route toward later-generation sensing rather than a
permanent EEG dependency. If peripheral signals can reproduce the decisions and
outcomes required for effective closed-loop TMR, the rationale for retaining
EEG would weaken. Conversely, systematic disagreement in physiologically
important intervals would identify functions for which direct
electrophysiological observation remains necessary.


\subsection{Generalization, Longitudinal Evaluation, and Responsible Translation}

Demonstration of a reference-system effect would establish neither broad
cognitive enhancement nor readiness for deployment. The final research phase
must determine whether the intervention survives changes in task, population,
time scale, and operating environment without losing efficacy,
interpretability, or acceptable effects on sleep.

\noindent\textbf{Generalization across memory domains.}

TMR effects should be evaluated separately across memory systems and learning
tasks rather than generalized from a single successful paradigm. Verbal paired
associates, spatial learning, procedural skills, emotional memories, and more
complex educational material differ in encoding structure, neural dependence,
cue representation, and behavioral outcome. The effect observed in one domain
should therefore be treated as evidence for that experimental class until
cross-domain replication demonstrates broader applicability.

Generalization studies should preserve task-appropriate controls and report both
absolute learning performance and the differential effect of cueing. Transfer
to uncued but related knowledge should be measured explicitly rather than
assumed. Evidence that TMR strengthens one set of cued items does not by itself
support claims of general improvement in learning ability, intelligence, or
memory capacity.

\noindent\textbf{Population heterogeneity and individual response.}

Future studies should progressively include participants differing in age,
baseline memory performance, sleep architecture, sensory threshold, and other
characteristics likely to influence intervention response. Expansion to
clinical populations should occur only when the intended therapeutic question
is independently justified and when the physiological and behavioral outcome
measures are appropriate to that population.

Group-level efficacy should not be treated as evidence of predictable
individual benefit. Future trials should characterize the distribution of
responses, including reproducible non-response and possible adverse-response
patterns, and should distinguish stable individual differences from
night-to-night variability.

\noindent\textbf{Longitudinal efficacy and sleep integrity.}

Repeated-night studies are necessary to determine whether TMR effects persist,
accumulate, diminish, or vary across extended use. Relevant outcomes include
retention over longer intervals, night-to-night effect stability, cue
habituation, changes in sensory responsiveness, sleep continuity, stage
distribution, arousal burden, and effects on uncued material.

Potential outcomes such as habituation, progressive sleep fragmentation, or
interference with other memory processes should be treated as hypotheses to be
measured rather than consequences assumed to occur. Longitudinal validation
should therefore evaluate both efficacy and the preservation of normal sleep
physiology under repeated exposure.

A sustained intervention should be considered scientifically supported only
when behavioral benefit remains reproducible without evidence that the
procedure progressively compromises the physiological state on which the
intervention depends.

\noindent\textbf{Responsible translation and evidentiary claims.}

The ethical problem surrounding closed-loop TMR extends beyond data protection
because the system intervenes in memory processing during a state in which the
user cannot actively supervise each decision. Future studies should therefore
address not only informed consent and protection of physiological data, but also
control over which memories are eligible for cueing, transparency regarding
automated intervention, the ability to suspend or withdraw participation, and
the distinction between experimentally demonstrated effects and speculative
enhancement claims.

EEG recordings, sleep profiles, learning performance, and intervention histories
should be collected according to data-minimization principles and protected
according to the sensitivity and intended use of the information. Expansion
toward educational, occupational, clinical, or enhancement contexts should not
be treated as equivalent because the acceptable evidence, risk, and oversight
requirements differ across those applications.

Communication of efficacy should remain proportional to the evidence. A modest
average TMR effect should not be represented as guaranteed improvement, and
successful performance in a constrained task should not be advertised as
general memory enhancement. Particular caution is required before research is
extended to populations with reduced capacity to consent or to applications in
which participation could become coercive.

The long-term research objective is therefore not simply to produce a smaller,
more adaptive, or more scalable TMR system. It is to determine the minimal
physiological information and intervention complexity required to produce a
reproducible memory effect while preserving sleep and maintaining an
interpretable relationship between measurement, stimulation, and outcome.
Progression beyond the first-generation reference pathway should occur only
when an alternative demonstrates a measurable improvement in efficacy,
reliability, safety, or measurement burden without introducing greater
unresolved uncertainty.